%% ma: L11-B16-0006
%% lop: 11
%% chuong: 5
%% bai: 16
%% dang: 11.16.3
%% loai: TL
%% muc: VD
%% dap_an: "-"
%% nguon: "tai-lieu-hoc-tap-mon-toan-lop-11-theo-chuong-trinh-giao-duc-pho-thong-2018.pdf (Lớp Toán Thầy Duy, Phạm Lê Duy, THPT Chu Văn An), Chương 3 Bài 2, Dạng 2, Ví dụ 6 (tr.354)"
%% kien_thuc: "Lớp 11 Bài 16 (giới hạn của hàm số)"
%% trang_thai: cho-duyet
%% ly_do: "Dạng mới đề xuất (11.16.x), chờ giáo viên duyệt dạng; tài liệu gốc không có lời giải mẫu, đáp án do kho tự giải và kiểm bằng sympy."
\begin{ex}
Tính các giới hạn sau.

a) $\lim\limits_{x\to-\infty}\dfrac{-6x^5+7x^3-4x+3}{8x^5-5x^4+2x^2-1}$;

b) \[\lim\limits_{x\to+\infty}\dfrac{\sqrt{9x^2+2}-\sqrt[3]{6x^2+5}}{\sqrt[4]{16x^4+3}-\sqrt[5]{8x^4+7}};\]

c) $\lim\limits_{x\to-\infty}\dfrac{\sqrt{x^2+2x}+3x}{\sqrt{4x^2+1}-x+2}$;

d) $\lim\limits_{x\to-\infty}\dfrac{(2x-3)^{20}(3x+2)^{30}}{(2x+1)^{50}}$.
\loigiai{\sloppy a) Chia cho $x^5$: giới hạn $=\dfrac{-6}{8}=-\dfrac34$.

b) Chia tử và mẫu cho $x$ ($x>0$): tử $\to\sqrt9=3$ (vì $\dfrac{\sqrt[3]{6x^2+5}}x\to0$), mẫu $\to\sqrt[4]{16}=2$ (vì $\dfrac{\sqrt[5]{8x^4+7}}x\to0$). Giới hạn $=\dfrac32$.

c) Với $x<0$: $\sqrt{x^2+2x}=-x\sqrt{1+\frac2x}$, $\sqrt{4x^2+1}=-x\sqrt{4+\frac1{x^2}}$. Chia tử và mẫu cho $x$: tử $\to-1+3=2$; mẫu $\to-2-1=-3$. Giới hạn $=-\dfrac23$.

d) Chia tử và mẫu cho $x^{50}$: \[\dfrac{\left(2-\frac3x\right)^{20}\left(3+\frac2x\right)^{30}}{\left(2+\frac1x\right)^{50}}\to\dfrac{2^{20}\cdot3^{30}}{2^{50}}=\left(\dfrac32\right)^{30}.\]}
\end{ex}
